\documentclass[11pt,a4paper]{article}
\usepackage[utf8]{inputenc}
\usepackage{amsmath,amsfonts,amssymb}
\usepackage{geometry}
\geometry{margin=1in}

\title{\textbf{Technical Writing: Scientific Typesetting with \LaTeX}}
\author{Jonathan Jolly \\ Department of Computer Science and Engineering \\ MNIT Jaipur}
\date{\today}

\begin{document}
\maketitle

\section{Introduction}
\LaTeX{} is a document preparation system widely used in academia and technical fields. It separates document formatting from content creation, allowing researchers and engineers to focus on clear writing.

\section{Mathematical Typesetting}
Technical documents frequently require precise mathematical expressions. For example, the quadratic formula is represented as:
\begin{equation}
x = \frac{-b \pm \sqrt{b^2 - 4ac}}{2a}
\end{equation}

Inline mathematical equations such as $\int_{0}^{\infty} e^{-x^2} dx = \frac{\sqrt{\pi}}{2}$ format cleanly alongside body text.

\section{Key Advantages}
\begin{itemize}
    \item Deterministic compilation across all operating systems.
    \item Robust automated referencing and citation management via Bib\TeX.
    \item Professional typographical quality.
\end{itemize}

\end{document}
